% answers_short.tex —— 2006 年·全国卷Ⅰ·理科 · 简短答案速查
% 由 tools/make_short_answers.py 自 answers.tex 的 \gksolhead/\gkanswer 提取，
% 并与 2006/2006_数学真题_全国1理_答案要点.md 逐题比对后生成，勿手改（改源 MD 后重跑）。

\begin{gknotice}
\textbf{使用说明}：本节只给最终答案便于快速核对，完整推导见解析版。答案为编写者独立推导，\textbf{非官方标准答案}。
\end{gknotice}

\gkpart{选择题}
\begin{center}\begin{tabular}{ccccccccccccc}
\toprule
题号 & 1 & 2 & 3 & 4 & 5 & 6 & 7 & 8 & 9 & 10 & 11 & 12 \\
\midrule
答案 & B & D & A & B & C & B & C & A & D & B & B & B \\
\bottomrule
\end{tabular}\end{center}

\gkpart{填空题}
\begin{itemize}[leftmargin=2.6em]
  \item[\textbf{13.}] $\dfrac\pi3$
  \item[\textbf{14.}] $11$
  \item[\textbf{15.}] $2400$
  \item[\textbf{16.}] $\dfrac\pi6$
\end{itemize}

\gkpart{解答题}
\begin{itemize}[leftmargin=2.6em]
  \item[\textbf{17.}] $A=\dfrac\pi3$ 时取得最大值，最大值为 $\dfrac32$。
  \item[\textbf{18.}] (1) $P(\xi=0)=\dfrac{125}{729}$，$P(\xi=1)=\dfrac{100}{243}$，$P(\xi=2)=\dfrac{80}{243}$，$P(\xi=3)=\dfrac{64}{729}$；(2) $E\xi=\dfrac43$
  \item[\textbf{19.}] (1) 见解析版；(2) $\sqrt{\dfrac23}$
  \item[\textbf{20.}] (1) 点 $M$ 的轨迹为 $\dfrac4{x^2}+\dfrac1{y^2}=1$ 的第一象限分支，且 $x>2,y>1$；(2) $|\overrightarrow{OM}|_{\min}=3$
  \item[\textbf{21.}] (1) $0<a\leqslant2$ 时在定义域的两个区间上递增；$a>2$ 时在 $(-\infty,-r)$、$(r,1)$、$(1,+\infty)$ 上递增，在 $(-r,r)$ 上递减，其中 $r=\sqrt{1-\dfrac2a}$；(2) $0<a\leqslant2$
  \item[\textbf{22.}] (1) $a_1=2$，$a_n=4^n-2^n$；(2) $\displaystyle\sum_{i=1}^nT_i=\dfrac32\left(1-\dfrac1{2^{n+1}-1}\right)<\dfrac32$
\end{itemize}
